\section{CLR: a VM que fez generics de outro jeito}
\label{sec:hist-clr}

\paragraph{O padrão.}
A plataforma .NET é especificada pelo padrão
\citetitle{doc:ecma335} (6ª edição, junho de 2012, equivalente ao
ISO/IEC~23271). A página da Ecma descreve a
\termo{CLI}{\emph{Common Language Infrastructure}: o padrão que define
como programas em várias linguagens rodam sobre uma mesma VM.} como uma
infraestrutura para que programas em várias linguagens rodem em vários
ambientes sem serem reescritos, e lista as partes do padrão: a
\emph{Virtual Execution System} (VES, a máquina virtual), os metadados
e o seu formato de arquivo, e o conjunto de instruções
\emph{Common Intermediate Language} (CIL). O CLR é a implementação da
Microsoft. Troque os nomes e a JVM aparece: a especificação da JVM
(\citetitle{doc:jvms25}) define a máquina, o formato \texttt{.class} e o
conjunto de instruções.

\paragraph{A diferença que importa: generics.}
Em Java, generics são resolvidos pelo compilador e apagados: a JLS chama
isso de \emph{type erasure}, \enquote{a mapping from types (possibly
including parameterized types and type variables) to types (that are
never parameterized types or type variables)}
(\citetitle{doc:jls25}, §4.6). Em runtime, \texttt{List<String>} e
\texttt{List<Integer>} são a mesma classe.

O CLR fez outra escolha. \textcite{paper:clr-generics} descrevem generics
com \enquote{exact run-time types}: a VM sabe, em runtime, que uma lista
é de inteiros. A implementação aproveita o JIT para fazer
\enquote{just-in-time type specialization}, com código compartilhado
entre instanciações de mesma representação; os primeiros resultados
sugerem que quem usa generics não paga sobrecusto por isso.

\paragraph{O que isso explica na JVM.}
\begin{itemize}
\item Por que não existe \texttt{new T[]} nem \texttt{instanceof
  List<String>} em Java: a informação foi apagada antes de chegar à VM.
\item Por que \texttt{List<Integer>} guarda objetos \texttt{Integer} e
  não \texttt{int}: com erasure, a VM só vê \texttt{Object}, e o
  autoboxing da seção~\ref{sec:autoboxing} é o preço. Uma VM que conhece
  o tipo exato pode especializar o código para o tipo primitivo, como o
  paper do CLR descreve.
\item Que erasure é uma escolha da plataforma Java, não um limite de
  toda VM com bytecode: o CLR mostra a alternativa, com o custo de levar
  os tipos genéricos para dentro da VM e do JIT.
\end{itemize}
